\begin{claim}{AC.40}{Proved: semilinear simultaneous regularity from all relative tests}
Let \(G\) be a finite solvable group of order prime to \(p\),
let \(V\) be a finite-dimensional \(\mathbf F_p\)-vector space
with a \(G\)-action, and let \(F\in\operatorname{GL}(V)\) and
\(\phi\in\operatorname{Aut}(G)\) satisfy
\[
                         F(gv)=\phi(g)F(v).
\]
Assume that for every positive integer \(r\) and every
\(\phi^r\)-stable pair of subgroups \(A\triangleleft B\)
with \(B\triangleleft G\) and \(B/A\) cyclic of prime order, one has
\[
                  (V^A)^{F^r}\supsetneq(V^B)^{F^r}.
                  \tag{AC.40.1}
\]
Then there is \(v\in V^F\) such that
\[
                    \operatorname{Stab}_G(e_Hv)=H
                    \quad\text{for every normal }H\triangleleft G.
                    \tag{AC.40.2}
\]
The subgroups \(H\) in this conclusion need not be
\(\phi\)-stable.

For a pointed geometrically Galois odd prime-to-\(p\)
etale cover \(D/k\to C/k\), take
\(V=J_D[p](\overline k)\), with its actual arithmetic
Frobenius. Failure of every relative cyclic-prime
class-number test over intermediate pointed bases which are
geometrically normal over \(C\), and over every finite
constant extension, implies (AC.40.1).
Consequently it supplies the class in (AC.40.2).
The rational class need not be a point difference.
\end{claim}
\begin{proof}
We induct on \(|G|\). Choose a minimal nontrivial subgroup
\(N\subseteq G\) which is normal in \(G\) and stable under
\(\phi\). It is elementary abelian of order a power of
some prime \(\ell\ne p\). Indeed its derived subgroup,
its Sylow subgroups, and its subgroup of \(\ell\)-th
powers would otherwise give smaller characteristic
subgroups; solvability rules out the nonabelian
characteristically simple case.

The averaging projector gives a \(G,F\)-stable splitting
\[
                         V=U\oplus W,\qquad
                         U=V^N,\quad W=(1-e_N)V.
\]
The hypothesis descends to the \(G/N\)-module \(U\),
so induction gives \(u\in U^F\) satisfying (AC.40.2)
for that quotient. For every normal \(H\) in \(G\),
\[
                         \operatorname{Stab}_G(e_Hu)=HN.
                         \tag{AC.40.3}
\]

Decompose \(W\) by the commuting projectors \(e_H\),
for all normal subgroups \(H\) of \(G\), and by the
simple factors of the commutative semisimple algebra
\(\mathbf F_p[N]\). The former joint blocks are
indexed by kernels \(L\triangleleft G\) of irreducible
\(G\)-constituents: on a block with kernel \(L\),
\[
                         e_H=1\quad\Longleftrightarrow\quad H\subseteq L.
\]
This follows from Maschke's theorem, since for normal
\(H\) the \(H\)-invariants of an irreducible constituent
are either the whole constituent or zero. We group
together all constituents with the same kernel; no
choice of a splitting field is needed for that grouping.
The further \(N\)-character factors correspond to
\(p\)-power orbits of nontrivial characters of \(N\).
All characters in such a factor have the same kernel.

The operator \(F\) permutes these refined blocks.
For each orbit of blocks whose direct sum has a
nonzero \(F\)-fixed space, choose a nonzero fixed
vector on that orbit, and add these choices to obtain
\(w\in W^F\). Each chosen vector has a nonzero component
in every block of its orbit. On an \(N\)-character
factor, a nontrivial character value minus one is
invertible; thus every nonzero component has exactly
the corresponding character kernel as its stabilizer
in \(N\). Consequently, simultaneously for all normal
\(H\), our \(w\) realizes the pointwise kernel:
\[
 \operatorname{Stab}_N(e_Hw)=
 \ker\bigl(N\text{ acting on the set }e_H(W^F)\bigr).
                         \tag{AC.40.4}
\]
Here \(e_H(W^F)\) need not be a \(G\)-submodule or
even equal \((W^H)^{F^r}\) for a field of definition
of \(H\). We do not make either of these replacements.

We prove that the kernel in (AC.40.4) is \(H\cap N\).
Suppose that it contains \(n\in N\setminus H\).
If \(n\) acted trivially on \(V^H\), so would all its
\(G\)-conjugates, because \(V^H\) is \(G\)-stable. Put
\[
 B=H\langle gng^{-1}:g\in G\rangle.
\]
Then \(B\triangleleft G\), \(B/H\) is nontrivial
elementary abelian of exponent \(\ell\), and
\(V^H=V^B\). Choose a hyperplane \(A/H\) in \(B/H\).
It gives \(A\triangleleft B\), \(B/A\simeq C_\ell\),
and \(V^A=V^B\), contradicting (AC.40.1) after a
common power making \(F\) and \(\phi\) the identity.
Thus \(n\) does not act trivially on \(V^H\).
Since \(N\) acts trivially on \(U\),
there is an \(N\)-character \(\chi\) in \(W^H\)
with \(\chi(n)\ne1\).

Among all normal constituent kernels \(L\supseteq H\)
supporting a nonzero character in the line
\[
                 \{\chi^a: a\in\mathbf F_\ell^\times\},
\]
choose one maximal under inclusion. Put
\(K=\ker\chi\subset N\). A surviving constituent with
kernel \(L\) shows \(L\cap N\subseteq K\), so
\[
                         LK\triangleleft LN,\qquad
                         LN/LK\simeq C_\ell.
\]
Modulo \(L\), the first subgroup is
\(K/(L\cap N)\) inside the abelian group
\(N/(L\cap N)\), which proves the asserted normality.
Consider the relative module
\[
                         V^{LK}/V^{LN}.
                         \tag{AC.40.5}
\]
Its constituents have a normal \(G\)-kernel containing
\(L\), and a nontrivial \(N\)-character in the displayed
line. Maximality therefore forces their \(G\)-kernel
to be exactly \(L\). Thus (AC.40.5) consists precisely
of the \(L\)-kernel blocks with nontrivial \(N\)-character
kernel \(K\).

No full \(F\)-orbit meeting this relative module can
have an \(F\)-fixed vector. If it did, our choice of
\(w\) would put a nonzero component in its selected
block. Since \(H\subseteq L\), that component survives
\(e_H\); every nonzero multiple of \(\chi\) detects
\(n\). This contradicts \(n\in\operatorname{Stab}_N(e_Hw)\).

Let \(r\) be the least positive integer for which
\(\phi^r(L)=L\) and \(\phi^r(K)=K\). The pair \((L,K)\)
is determined by each refined block in (AC.40.5).
Hence \(r\) divides the length \(t\) of its full
\(F\)-orbit of blocks. The part of that orbit appearing
in (AC.40.5) is exactly one cycle under \(F^r\).
A fixed vector of \(F^r\) on this cycle is equivalent
to a fixed vector of \(F^t\) on one block, and hence
to a fixed vector of \(F\) on the full orbit.
There is none by the preceding paragraph.

The averaging splitting of (AC.40.5) is \(F^r\)-stable,
so we conclude
\[
                         (V^{LK})^{F^r}=(V^{LN})^{F^r}.
\]
This contradicts (AC.40.1) for the cyclic-prime pair
\(LK\triangleleft LN\). Therefore
\(\operatorname{Stab}_N(e_Hw)=H\cap N\) for every normal
\(H\). Combining this with (AC.40.3) in the
\(G\)-stable direct sum \(U\oplus W\) proves that
\(v=u+w\) satisfies (AC.40.2), exactly as in AC.39.

For the geometric application, odd-order groups are solvable
by the verified Feit--Thompson theorem used in AC.3.
A \(\phi^r\)-stable
subgroup defines its pointed intermediate curve over
\(k_r\). Prime-to-\(p\) norm and pullback identify
the relevant \(p\)-torsion invariants with
\(V^A,V^B\); their strictness is exactly the presence
of rational new \(p\)-torsion, hence failure of (R)
for \(D/A\to D/B\) over \(k_r\). The deck group of
this relative cyclic-prime cover may be nonconstant.
If (R) succeeded, AC.20 would make that cyclic group
constant after its exact Frobenius orbit length,
raising the relative class number to that power and
preserving its \(p\)-unit property. Thus failure of
all constant-deck relative tests over every extension
indeed implies the stated hypothesis.
\end{proof}

\paragraph{The normal-level distinction remains necessary.}
The conclusion concerns all normal quotients of \(D/C\).
The relative cyclic cover isolated in the proof can
have source \(D/(LK)\) which is not normal over \(C\),
although its base \(D/(LN)\) is normal over \(C\).
Accordingly the compactness consequence for the
normal-level ghost module is a one-way implication
from hereditary failure of all relative tests.
The reverse assertion needs freeness at the nonnormal
source as well; it does not follow merely from
regularity at every normal quotient.
